% Nouveau EGA V (cote 25): archive cote 25, site page 53 (PDF page 54).
% typescript. Draft compared: none.
% AI-transcribed by Claude Opus 5.5 (claude-opus-5-5, adaptive thinking, effort high), Message Batches API, batch ega25b (msgbatch_01A7n5YuXubUPNK1zsepi8RG), 2026-11-13; original AI transcription, subsequently checked against both supplied scan parts; see PAGES_CHECKED.tsv and CHECK_DECISIONS.tsv.
\scanpage{25/53}
\margin{\hand{\unsure{46|48}}}

$\overset{\hand{\vee}}{P}$ (correspondant à la question d'étudier les sections hypérplanes de X
passant par un point donné x, ou tangent à X en un point \deleted{simple}\inserted{lisse} donné),
form\over{e}{é} donc des hyperplans contenant une sous-variété linéaire $L^{o}$ de P
(resp. un point,  ou l'image d'un espace tangent de X en un point lisse,
dans les deux cas envisagés), et on se pose la question de savoir si
pour le point générique \hand{$\xi$} de L  (donc pour tous les points d'un ouvert
non vide de L),  $Y_{\hand{\xi}}$  est \xout{géom intègre ou du moins} géom irréd., de
dim  dim X -1.   C'est là une variante  de Bertini, qui devrait figurer
dans  le n°\hand{3}  et se traite par exactement la même méthode\inserted[typed]{ou, si on veut, s'y ramène.}  La question
dimensionnelle s'exprime simplement par  $f(X)\hand{\not\subset} L^{o}$ , \xout{Il y a donc un}
i.e.  $X' =$ \xout{X=} $f^{-1}(P - L^{o})$   est un ouvert dense de X, \xout{et la question}
Soit \unsure{Q|N} l'espace projectif des hypérplans passant par $L^{o}$  (NB si  $L^{o}$ est
défini par un sous-espace vectoriel $F^{o}$ de E,  on a   $Q = P($\deleted{E/}$F^{o})$ ) et
considérons le morphisme  canonique  (déduit de \xout{E\illegible{1}E/$F^{o}$} $F^{o}\hand{\longrightarrow} E$ , cf
Chap II)
\[ u :\ P - L^{o} \hand{\longrightarrow} Q, \]
et considérons
\[ g = u f' :\quad f^{-1}(P - L^{o}) = X' \hand{\longrightarrow} Q , \]
alors \xout{la famille des X' \illegible{1} n'est autr}  $L \hand{\simeq} \overset{\hand{\vee}}{Q}$ ,  et la famille des
$X'_{\hand{\xi}}$  $(\hand{\xi}\hand{\in} L)$  n'est autre que la famille des ``sections hyperplanes'' rela-%
tif à ce morphisme  g . D'autre part, on voit aussitôt que pour tout
$\hand{\xi}\hand{\in} L$  ``général''  $X'_{\hand{\xi}}$ est dense dans $X_{\hand{\xi}}$ , donc  $X'_{\hand{\xi}}$  est géom irréd sss
X  l'est.  Ceci posé,  le théorè\unsure{m|p}e de Bertini-Zariski nous montre que
nous avons la conclusion voulue \xout{dix} d'irréductibilité pourvu que
dim $g(X') \hand{\geqslant} 2$. (A vrai dire, tant qu'à faire on pourrait signaler une
réciproque dans \hand{3}.1. , savoir : \xout{Y est gé} si X est géom irréd, Y  est
géom irréd sss  ou bien  dim $f(X)$; $\hand{\neq}$ \unsure{2|1}, ou bien  dim $f(X)=1$ , et $f(X)$
est contenu dans une droite\hand{D} (définie sur $\overset{\hand{\rule{0.6em}{0.35pt}}}{k}$ )\inserted{et la fibre générique de $X\to D$ est géom. \unsure{irréd}}\hand{)}; cela permet également

\margin{\hand{53}}

% NOTE: Beside dim dim X-1 the left margin has several parallel pencil strokes; no textual slash is supplied.
